%! Planning a Study — Unit 3, Lesson 3.1 — Warm-Up (Answer Key)
% Mirrors the blank exactly apart from the package swap, the pageheader, and the answers.
% NO teachernote here — teacher prose lives in the lesson plan.
% ONE PAGE, blank and key. Three items; do not add a fourth.
% GRADUAL RELEASE: the warm-up activates prior knowledge before the guided notes.
%   Item 1 — a news site's later-start poll, 2,400 responses, 85% yes: who answers a poll like
%            this? Seeds Row 1 (Ask how) — the room's answer becomes "non-chance method".
%   Item 2 — 900 students and a lunch question: ask all of them or some? Seeds Row 2 (the
%            population, a sample, a census) — the trade-off is time against completeness.
%   Item 3 — spiral to Unit 2: explanatory and response variables, and why an observed
%            association is not a cause. Seeds Row 3 (observational study vs experiment).
% Every display is given; nothing is sketched.
\documentclass[12pt]{article}
\usepackage{apstats-article}
\usepackage{apstats-key}   % pulls in -boxes; do NOT also load -boxes

\begin{document}

\pageheader{Unit 3, Lesson 3.1}{Warm-Up --- Answer Key}
% NAMESTRIP: no \namedateperiod here — the name row lives on the cover only.

\vspace{0.28in}

\begin{enumerate}[leftmargin=1.9em, itemsep=0.42in]

  \item A news site posts an online poll: \emph{``Should the school day start later?''} Of the
        2{,}400 students who respond, \textbf{85\%} say yes. Who is most likely to answer a poll
        like this one? Does 85\% describe \emph{all} students?

        \vspace{12pt}
        \ansline{Students who care about the question --- mostly those who want a later}\\[18pt]\ansline{start. No: they chose to answer, so 85\% likely overstates support}\\[18pt]\ansline{among all students.}

  \item A school of 900 students wants to know which new lunch option students prefer. Give one
        reason to ask \emph{all} 900 students, and one reason to ask only some of them.

        \vspace{12pt}
        All 900: \ansline{no one is left out; we learn exactly what everyone thinks.}\\[18pt]
        Only some: \ansline{asking 900 takes too much time and effort.}

  \item \textbf{From Unit 2.} In a study of 200 students, those who slept more tended to score
        higher on a math test.

        \vspace{12pt}
        Explanatory variable: \ans{Hours of sleep} \qquad Response variable: \ans{Test score}
        \vspace{16pt}

        Does this show that more sleep \emph{causes} higher scores? Why or why not?
        \vspace{12pt}

        \ansline{No. The researchers only observed; something else (study habits, stress)}\\[18pt]\ansline{could drive both sleep and scores. Association is not cause.}

\end{enumerate}

\end{document}
